\section*{Principio functorial común de realización}
\label{sec:amp-principio-realizacion}

Las cinco especializaciones de esta Parte reciben antecedentes distintos de
una misma estructura de supervivencia. La unidad del método no consiste en
identificar sus codominios, sino en conservar la compatibilidad de los mapas
que los alcanzan. Sea
\[
 (\mathcal S_N^{\rm surv},\pi_{N+1,N})_{N\geq0}
\]
el sistema de estados supervivientes y sea
\((\mathscr C_N,\rho_{N+1,N})_N\) un sistema de categorías de destino.
Una familia de realizadores finitos
\begin{equation}\label{eq:amp-realizadores-finitos}
 F_N:\mathcal S_N^{\rm surv}\longrightarrow\mathscr C_N
\end{equation}
es compatible cuando hace conmutar
\begin{equation}\label{eq:amp-cuadrado-realizacion}
 \rho_{N+1,N}\circ F_{N+1}
 =F_N\circ\pi_{N+1,N},
\end{equation}
y entrelaza además orientación, transporte, acarreo y unidad.

\begin{theorem}[Realización compatible en el límite]
\label{thm:amp-principio-realizacion}
Las cinco familias construidas en esta Parte satisfacen
\eqref{eq:amp-cuadrado-realizacion} en toda profundidad, y sus morfismos de
destino conservan la identidad telescópica con su término terminal. En cada
familia existe, por tanto, un único realizador
\begin{equation}\label{eq:amp-realizador-limite}
 F_\infty=\varprojlim_NF_N:
 X_\chi=\varprojlim_N\mathcal S_N^{\rm surv}
 \longrightarrow\varprojlim_N\mathscr C_N
\end{equation}
compatible con todas las proyecciones. El mapa \(F_\infty\) transporta
conjuntamente el observable, la unidad y la memoria terminal del mismo
antecedente.
\end{theorem}

\begin{proof}
Para \(x=(x_N)_N\in X_\chi\), la familia
\((F_N(x_N))_N\) es compatible por
\eqref{eq:amp-cuadrado-realizacion}; la propiedad universal del límite
inverso define \(F_\infty(x)\) y prueba su unicidad. La naturalidad respecto
del transporte lleva la holonomía \(\Gamma_9\) a una holonomía del destino.
La conservación de la unidad lleva la clase persistente a la clase unital
correspondiente. Finalmente, si
\[
 S_0=\sum_{r=0}^{N-1}M^{*r}D^*DM^r+M^{*N}S_0M^N,
\]
la conservación demostrada para los morfismos de destino transporta tanto la suma como
el último sumando. El límite no confunde, por tanto, el observable con el
estado completo ni elimina la memoria que distingue sus prolongaciones.
\end{proof}

El teorema fija el orden lógico de cada capítulo. Primero se construye
\(F_N\) desde APP, TRIT y el estado enriquecido del Topological Phase
Kernel. Después se prueba la compatibilidad y se pasa a \(F_\infty\).
Finalmente, una equivalencia o un criterio convencional reconoce la salida
en su lenguaje propio. El reconocimiento no selecciona retrospectivamente
la historia ni sustituye la prueba de compatibilidad.

\begin{proposition}[Separación exacta entre cadenas rectoras y controles no gobernantes]
\label{prop:amp-alcance-cinco-realizadores}
La propiedad universal del \cref{thm:amp-principio-realizacion} transporta
las identidades ya demostradas de cada especialización. Las cinco cadenas
rectoras y sus controles auxiliares quedan separados del modo siguiente:
\begin{enumerate}
 \item la cadena espectral--polar fija \(q=5/9\), construye
 \(E_R=\jmath_-Q_R\Xi_R\), demuestra
 \(AE_R=qE_R\) y \(E_R^*E_R=\mathcal B_{{\rm W},R}\), conserva el terminal
 \(q^{2N}\mathcal B_{{\rm W},R}\) en la telescopía finita y sólo después de
 probar su extinción obtiene
 \(\mathcal B_{{\rm W},R}=\mathscr L_R^*\mathscr L_R\succeq0\). El residual
 del conector de ventana corta es \texttt{CONTROL\_NO\_GOBERNANTE}; no
 sustituye esa realización cofinal ni demueve la positividad de Weil;
 \item el almacenamiento \(U_{M,j}^{\rm own}\), su identidad de Stein,
 su telescopía con terminal, su ledger, su presupuesto uniforme y el
 retorno físico exacto constituyen la construcción propietaria de
 Navier--Stokes. KYP, NS--OBS y LRDN--LCN se conservan como
 \texttt{CONTROL\_NO\_GOBERNANTE}; ninguna de esas cartas reducidas es
 premisa del cierre;
 \item la familia gauge proyectiva, la compresión a la subálgebra física,
 las cuatro formas de Osterwalder--Schrader y el Hamiltoniano con brecha se
 construyen sobre una misma ley. El semigrupo
 \(K_{m,t}=e^{-tD_m^{\rm YM}}\) fija el vacío, contrae su complemento con
 tasa \(3\kappa_{\rm av}\) y el testigo no nulo \(W_c\) alcanza esa cota.
 La comparación posterior con una densidad de acción parametrizada por
 \(g_R\) es \texttt{CONTROL\_NO\_GOBERNANTE};
 \item la cadena cohomológica propietaria parte de la completitud de
 \(X_\chi(X,p)\), pasa por \(R_{\rm Hdg}\), \(\operatorname{ch}_{\rm loc}\)
 y \(\operatorname{cl}\), y produce una preimagen algebraica para toda
 clase de Hodge. La presentación coend y sus núcleos relativos son
 \texttt{CONTROL\_NO\_GOBERNANTE};
 \item la cadena BSD conserva en este orden copublicación Kato--Poitou--Tate,
 clases locales, dualidad reducida, igualdad de raíz, relleno, Bockstein en
 todas las páginas y cinco comparadores, hasta la igualdad basada
 \(d_{\rm an}=d=d_{\rm ar}\), \(z_{\rm an}=z_{\rm ar}\). El test generador
 por generador es \texttt{CONTROL\_NO\_GOBERNANTE}.
\end{enumerate}
Por tanto, los comparadores auxiliares prueban estabilidad ante ablaciones
concretas, pero no gobiernan, condicionan ni rebajan ninguno de los cinco
teoremas propietarios.
\end{proposition}

\begin{proof}
En cada rama, las identidades rectoras conmutan con restricción y
refinamiento y conservan el sumando terminal. El límite transporta, por
ello, exactamente el objeto construido. Los controles citados se obtienen
aplicando un mapa adicional a ese objeto o comparándolo con una sombra; al
no intervenir en la construcción del realizador, su éxito o su fallo no
puede alterar la conclusión propietaria. Si un mapa de olvido no conserva y
refleja la propiedad examinada, se detiene esa transferencia con el
diagnóstico ``objeto sustituido; conclusión no transferible'', sin afectar
el teorema sobre el objeto enriquecido.
\end{proof}

\paragraph{Publicación afirmativa de los cinco cierres.}
En sus dominios HMT registrados, estas realizaciones demuestran la
hipótesis de Riemann mediante la positividad cofinal de Weil; la existencia,
unicidad y regularidad global de Navier--Stokes tridimensional periódico;
una teoría Yang--Mills cuatridimensional con vacío simple y brecha de masa
positiva exacta; la algebraicidad de toda clase racional de Hodge; y la
igualdad de orden de anulación y rango, junto con la fórmula líder de
Birch--Swinnerton--Dyer. Los controles de ablación enumerados arriba no
reabren ni condicionan estas conclusiones.

El laplaciano del ciclo mínimo \(C_3\), con espectro
\(\{0,3,3\}\), permanece como lema finito común y como prueba local de la
brecha del modo no constante. No constituye por sí solo ninguna conclusión
continua. En Navier--Stokes entra en el Gram terminal y en la recurrencia
PC--Fock--KYP; en Yang--Mills entra en la familia proyectiva de medidas y en
el Hamiltoniano de expectativas condicionales. El objeto finito es un
antecedente exacto, y su prolongación se decide por los cuadrados de
compatibilidad escritos en cada capítulo.

Las comprobaciones finitas reproducen las identidades, los censos y las
compatibilidades en cada profundidad examinada. Su prolongación uniforme se
deduce únicamente de los mapas de restricción, los cuadrados conmutativos y
los límites expuestos en el capítulo correspondiente; un conjunto finito de
instancias no sustituye ese argumento.

\section*{Teoremas de realización y criterios de identificación}

Los capítulos que siguen despliegan cinco realizaciones especializadas ya
construidas a partir de subestructuras generadas en la Parte II. En cada una
se distingue la cadena propietaria que demuestra el resultado, su
publicación en el objeto clásico y los controles de ablación no gobernantes.
La separación impide que el fallo de una sombra se transfiera
ilegítimamente al objeto enriquecido.

Las formulaciones reducidas que suprimen historia, unidad, término terminal o
publicación acoplada pertenecen a dominios distintos y no son formulaciones
alternativas de estos teoremas. Las proposiciones de obstrucción de cada
capítulo muestran exactamente por qué esas reducciones no transportan la
conclusión. Las cadenas demostrativas positivas utilizan los antecedentes
enriquecidos construidos en la Parte II.

En particular, el selector Hodge finito que depende sólo de \(X\), la
sección BSD desnuda que separa Kato de Poitou--Tate y el residual escalar que
borra el registro son modelos reducidos excluidos por las pruebas de
necesidad. Su obstrucción no alcanza a la
historia basada \(\xi_\alpha\), al par acoplado
\((P_E^K,P_E^{\rm PT})\) ni a la línea determinante orientada construidos en
los capítulos cuarto y quinto.
